\documentclass[11pt,a4paper]{article}
\usepackage[utf8]{inputenc}
\usepackage[T1]{fontenc}
\usepackage{lmodern,amsmath,amssymb,textcomp}
\usepackage[margin=24mm]{geometry}
\usepackage{longtable,booktabs,array,enumitem}
\usepackage[hidelinks]{hyperref}
\setlength{\parindent}{0pt}
\setlength{\parskip}{6pt}
\emergencystretch=3em
\begin{document}
{\LARGE\bfseries An exact 61-line 1190-triangle construction: a priority-review manuscript\par}\medskip
{\small Provisional mathematical authors: Rohith Poola\par}\medskip
{\small Judging and evaluation record: the submissions discussed in this document have been judged and evaluated by Alejandro Zarzuelo Urdiales for OpenMath 2026. This dated review records the stated verification, classification and unresolved holds; it is a preliminary evaluation rather than a signed award decision.\par}

{\small Event provenance: OpenMath 2026 ran 27 September-2 October 2026 with worldwide online participation. Detailed organizer listings specify Harvard Innovation Labs for the 27 September in-person event; the OpenMath programme advertises a hybrid kickoff at MIT. Sundai identifies its Harvard/MIT community. Organizer sources: \url{https://rsihouse.ai/openmath} ; \url{https://luma.com/yzp9abvr} ; \url{https://www.sundai.club/events/boston/recursive-learning-hack-with-harvard-innovation-labs}\par}

{\small Rohith Poola  -  Autonomy and Intelligence Lab, Northeastern University.\par}

{\small Affiliations follow the recorded author declarations or public institutional/profile sources. Preferred author names, author order and publication affiliations await author review. This editorial compilation does not establish anthology coauthorship.\par}

{\small Candidate manuscript  -  originality unresolved | OpenMath 2026 | 5 October 2026\par}\medskip
{\small Central source and adjudication archive: \url{https://github.com/alejandrozu/openmath-2026-judging}\par}

\subsection{Abstract}
The retained exact line data give a simple 61-line arrangement with 1190 bounded triangular faces. This exceeds the recorded local all-order baseline 1181, but historical priority for ordinary straight lines is unresolved. The manuscript remains outside the original-results anthology until that question is settled.

The supplied 61-line ordinary straight arrangement has 1190 bounded triangular faces; historical novelty remains unresolved.

\subsection{The surviving contribution from the catalogue}
The catalogue contains 66 arrangements, many originally marked new because a gallery omitted their orders. A missing gallery value does not establish a frontier result. Against stronger recorded pre-event general and finite bounds, eleven of the twelve proposed surviving rows are tied or dominated. The 61-line 1190 construction is the remaining candidate.

Its data independently check exactly: no parallel pairs, no duplicate lines and 1830 distinct pair intersections, with 1190 bounded triangular faces under the consecutive-edge criterion. Coefficients are within the declared construction restrictions. This establishes the geometry of the witness, not its historical originality.

\subsection{Why the purported 1199 prior is inconclusive}
The located Bartholdi-Blanc-Loisel corollary gives a generic doubling operation for pseudolines. Applying it to a 31-line 299 seed yields 61 pseudolines with 1199 triangles. That does not show that the doubled pseudoline arrangement can be realized by straight lines.

Their straight-line proposition has special coordinate and crossing hypotheses. A31-line straight seed must satisfy those hypotheses before the 1199 consequence can be used. The ordinary 31/299 gallery is not enough. The earlier Forge-Ramírez Alfonsín recursion also distinguishes projective pseudolines from specially constrained straight constructions. The inspected established straight-line sequences do not contain 61.

Consequently the candidate should not be rejected solely from a secondary table that conflates those objects. Equally, failure to locate a qualifying seed does not certify originality. A specialist should supply either a dated 61-line straight certificate or the dated special31-coordinate seed from which it follows.

\subsection{Publication gate}
The paper should expose the exact seed and modification used in the 1190 construction, credit the predecessor and gallery data, and identify what is genuinely changed. Attach the rational line list and exact triangle-face certificate. State K(61)\ensuremath{\geq}1190 as a verified construction theorem, while withholding best-known or first-record language.

Rohith's39/470 witness is not a new numerical result because the pre-event organizer packet already proves the same bound. The general upper-bound conjecture and unrealized pseudoline SAT outputs remain separate research questions. The 1190 paper is included as a private held draft for a concrete originality decision, not as an accepted novel contribution.

\subsection{Verification and reproducibility}
No submitted Lean project was located for this packet. The independent exact straight-line geometry recount is an ordinary certificate check, and is not represented as Lean replay. Originality remains separately under review.

The preserved contribution record and source locator below identify the selected certificate. No new endpoint statement is invented for this editorial manuscript.

{\small\textbf{Sources and attribution:} \href{https://github.com/Rohith18p/rsi-kobon-triangles/blob/f462d8e18aea2a458376c523c9b6c2980237071f/kobon-triangles/bases/n61_1190_base31.json}{1. Rohith18p/rsi-kobon-triangles / n61\_1190\_base31.json}; \href{https://arxiv.org/html/0706.0723v1}{2. arXiv source: 0706.0723v1}; \href{https://link.springer.com/content/pdf/10.1007/PL00009373.pdf}{3. link.springer.com / PL00009373.pdf}; \href{https://zegalur.github.io/line-order/gallery/kobon.html}{4. zegalur.github.io / kobon.html}; \href{https://github.com/alejandrozu/ulam-opdp-difficulty-atlas}{5. alejandrozu/ulam-opdp-difficulty-atlas}; \href{https://github.com/alejandrozu/openmath-2026-judging/blob/d0ed487f487fa7ec814ff705ab55249983fe8707/OpenMath-Judging/review/entries/E57-K61.txt}{6. alejandrozu/openmath-2026-judging / E57-K61.txt}; \href{https://github.com/alejandrozu/openmath-2026-judging}{Central source and adjudication archive}.\par}

\end{document}
